\section{Comparison with Menyuk, Chen and Lee}
\label{sec:MCL-dictionary}

The restricted shared-mode Hamiltonian and the broader wave construction
in \cite{MenyukChenLee1983Integrable} give two distinct comparisons
with the canonical equations used here.

\subsection{The restricted shared-mode family}

Write the shared mode as $p$ and the remaining pairs as $q_j,r_j$,
$1\leq j\leq R$. Complexify the conjugates as independent partners.
With brackets $\{p,\bar p\}=-i$ and the same convention for the
other pairs, the action--angle Hamiltonian in
\cite[equation (1.1)]{MenyukChenLee1983Integrable} is
\begin{equation}\label{eq:MCL-star-H}
 H_{\rm MCL}=\sum_{j=1}^{R}\left[
 \frac{\Delta_j}{2}(q_j\bar q_j-r_j\bar r_j)
 -\frac{\epsilon_j}{2}
 (p\bar q_jr_j+\bar p q_j\bar r_j)\right].
\end{equation}
Here $\Delta_j$ is the frequency mismatch and $\epsilon_j$ the
coupling. For the equal-coupling family $\epsilon_j=\epsilon$, set
\[
 u=p,\quad\bar u=\bar p,\qquad
 v_j=i\bar q_j,\quad\bar v_j=iq_j,\qquad
 w_j=r_j,\quad\bar w_j=\bar r_j,
 \qquad1\leq j\leq R.
\]
This is a complex canonical change: for example,
$\{i\bar q_j,iq_j\}=-i$. It gives
\begin{equation}\label{eq:MCL-star-map}
 H_{\rm MCL}=-\frac12\sum_{j=1}^{R}\Delta_j
 (v_j\bar v_j+w_j\bar w_j)
 +\frac{i\epsilon}{2}\sum_{j=1}^{R}
 (uv_jw_j+\bar u\bar v_j\bar w_j).
\end{equation}
Thus the restricted model has our interaction arrays with
$(L,M,N)=(1,1,R)$, up to the constant coupling and time normalization.
Its original physical real form becomes $\bar v_j=-v_j^*$,
$\bar u=u^*$ and $\bar w_j=w_j^*$ under this complex dictionary;
the transformed coupling is imaginary. This is a different convention
from imposing actual conjugation on all partner arrays.

Menyuk, Chen and Lee give a complete set of $2R+1$ independent
commuting integrals for this restricted family, including coincident
mismatches \cite[Section~III]{MenyukChenLee1983Integrable}.
Their Section~V assumes equal mismatches as well as equal couplings:
the shared power is elliptic, its phase is an elliptic integral of
the third kind, and the remaining modes are recovered by quadratures.

\subsection{The broader matrix wave construction}

Their equation (2.13), attributed in its reality-reduced form to
Ablowitz and Haberman, is
\begin{equation}\label{eq:MCL-matrix-wave}
 \alpha_{hi}\partial_x N_{hi}
 =\partial_t N_{hi}
 +\sum_{\substack{j=1\\j\ne h,i}}^{d}
 N_{hj}N_{ji}(\alpha_{ji}-\alpha_{hj}),
 \qquad1\leq h,i\leq d,\quad h\ne i.
\end{equation}
For unequal source spectral constants,
$\alpha_{hi}=(c_i-c_h)/(i(a_i-a_h))$; coincident $a_h=a_i$,
$c_h=c_i$ permit separate choices. Section~IV takes
$N_{hi}(x,t)=N_{hi}(t)e^{\Delta_{hi}x}$ with
$\Delta_{hi}=\Delta_{hj}+\Delta_{ji}$, producing linear rates and
quadratic mode interactions.

Here is the reduction to our arbitrary-size canonical system.
Use the $H_0$ modes in \cref{eq:canonical-X} and write
\[
 \widetilde{\mathsf X}=E+Z,\qquad
 E=\diag(-2\mathsf B_L,-2\mathsf B_M,\mathsf B_N).
\]
The matrix $Z$ has zero diagonal blocks. On the three groups of
indices choose symmetric coefficients
\[
 \alpha_{LM}=0,\qquad\alpha_{LN}=1,\qquad
 \alpha_{MN}=-1,
\]
and zero coefficients within each group. These follow from the
complexified source parameters $a=i,-i,0$ and $c=0,0,1$ on the
respective groups. Define $(T_0)_{hi}=\alpha_{hi}Z_{hi}$.
Our $H_0$ equations give $i\dot Z=[E+Z,T_0]$. Substituting
\begin{equation}\label{eq:MCL-matrix-map}
 N(x,t)=i e^{-ixE}Z(t)e^{ixE}
\end{equation}
into \cref{eq:MCL-matrix-wave} gives exactly this commutator:
at $x=0$, $\partial_xN_{hi}=(E_{hh}-E_{ii})Z_{hi}$ and the
quadratic sum is $-[Z,T_0]_{hi}$. The exponential rates
$\Delta_{hi}=-i(E_{hh}-E_{ii})$ obey the required consistency
condition. Equal coefficients to the third group preserve the zero
diagonal blocks. This comparison applies to all positive $L,M,N$;
\cref{eq:canonical-phase-gauge} gives the symmetric phase convention.

The broader wave framework therefore already permits these
phase-dependent equations. The independent phase parameters are a
constant-diagonal extension relative to the printed 1991 Hamiltonian,
as detailed in Appendix~\ref{sec:JRW-dictionary}. Their determination
from prescribed BA data is the role they play in the present paper.
The explicit complete-integral calculation in the MCL paper concerns
the restricted shared-mode family; the all-size canonical count used
here is stated in \cref{thm:canonical-independence} and proved in
Appendix~\ref{sec:integrability-proof}.
